%==============================================================================
% Closing the Variance Gap: Spectral PINN for Lévy-Driven SPDEs
% Hubert Lipski
%
% Version 3: includes §5.4 (per-component ablation on Burgers-VG, 5x5)
% and §6 (Spectral FNO on four problems, 8 trainings).
% Contributions list updated; limitations updated; §1.4.3 sentence added.
%
% Save as paper_fno.tex, put references.bib next to it, drop figures/
% folder next to it, then:
%     pdflatex paper_fno && bibtex paper_fno && pdflatex paper_fno && pdflatex paper_fno
%==============================================================================

\documentclass[review,12pt]{elsarticle}

\usepackage[utf8]{inputenc}
\usepackage[T1]{fontenc}
\usepackage{amsmath, amsthm, amssymb}
\usepackage{booktabs}
\usepackage{multirow}
\usepackage{graphicx}
\usepackage{hyperref}

\graphicspath{{./}{./paper_fno/}}

\newtheorem{remark}{Remark}
\newtheorem{proposition}{Proposition}

\journal{Journal of Computational Physics}

\begin{document}

\label{sec:spectral-fno}

The spectral loss introduced in Section~\ref{sec:method} is
architecture-agnostic: nothing in the construction
(eq.~\ref{eq:losses}) requires the predictor to be a coordinate-MLP
PINN. To stress-test this, we apply the same three loss terms to the
Fourier Neural Operator (FNO) of Li et al.~\cite{Li2021FNO} ---
the canonical operator-learning surrogate --- and evaluate the
resulting \emph{Spectral FNO} on five problems spanning the
SPDE landscape: the three benchmarks of Section~\ref{sec:setup}
plus two additional problems chosen to test the loss outside the
L\'evy-driven regime.

\subsection{Setup}
\label{subsec:fno-setup}

\paragraph{Architecture.}
The base architecture is a 4-layer 3D FNO with width 16 and
$(8, 8, 16)$ Fourier modes along $(x, y, t)$, identical to the
canonical configuration in~\cite{Li2021FNO}. The total parameter
count is approximately $8.4 \times 10^{6}$.

\paragraph{Loss.}
The vanilla FNO is trained with pointwise MSE only:
$\mathcal{L} = \mathcal{L}_{\mathrm{MSE}}$. The Spectral FNO adds
the three distributional terms with the same weights used in the
Spectral PINN:
\begin{equation}
\mathcal{L}_{\text{spec-FNO}} = \mathcal{L}_{\mathrm{MSE}}
+ \lambda_{xy} \mathcal{L}_{xy}
+ \lambda_{t}  \mathcal{L}_{t}
+ \lambda_{v}  \mathcal{L}_{\mathrm{var}},
\end{equation}
with $(\lambda_{xy}, \lambda_{t}, \lambda_{v}) = (1, 1, 50)$. The
spatial PSD weight is set to 1 (down from 3 in the PINN) because the
FNO already operates on Fourier coefficients architecturally; the
remaining weights are unchanged.

\paragraph{Problems.}
We evaluate on five problems:
(i) \emph{Burgers-VG} (Section~\ref{subsec:problems}),
(ii) \emph{Heston-VG} (Section~\ref{subsec:problems}),
(iii) \emph{FPL} --- the hierarchical 2D Fokker--Planck--L\'evy
system of Section~\ref{subsec:problems},
on a periodic domain at $\nu = 10^{-3}$ with weak random forcing,
solved by a pseudo-spectral integrator with 2/3-rule dealiasing,
and
(v) \emph{BS-L\'evy} --- a 2D Black--Scholes-type system tracking
two correlated Variance-Gamma jump-diffusion log-prices on a
$32 \times 32$ initial-spot grid.
The first three are L\'evy-driven; the last two add a
fluid-dynamics benchmark and a finance benchmark with shared
spatially-coupled VG drivers, respectively.

\paragraph{Training.}
Each (problem, variant) pair is trained on a fresh dataset of
20--32 trajectories (24 for FPL and BS-L\'evy, 20 for Heston-VG, 24 for Burgers-VG) at $32 \times 32 \times 128$ resolution with
4--8 validation trajectories. Optimisation is AdamW with cosine LR
$10^{-3} \to 10^{-5}$ and weight decay $10^{-5}$ for 50--60 epochs.
The complete sweep (4 problems $\times$ 2 variants = 8 trainings)
finishes in 1.34~h on a single RTX 4090. Test sets use independent base
seeds disjoint from training and validation.

\subsection{Results}
\label{subsec:fno-results}

\begin{table}[h]
\centering
\caption{Spectral FNO vs.\ vanilla FNO on four problems. Mean
$\pm$ standard deviation across test trajectories (8 each for
Burgers, BS-L\'evy; 12 for Heston-VG over two channels;
18 for FPL over three channels). Higher var.\ ratio is better
(target 1.0); lower log-PSD MSE is better. ``$\Delta$ var'' is
the absolute closure of the variance gap relative to the target.}
\label{tab:spectral-fno}
\begin{tabular}{l l c c c c}
\toprule
problem & variant & var.\ ratio & log-spat.\ PSD MSE & log-temp.\ PSD MSE & $\Delta$ var \\
\midrule
\multirow{2}{*}{Burgers-VG} &
  Vanilla FNO            & $0.761 \pm 0.040$ & $1.865 \pm 0.025$ & $1.218 \pm 0.014$ & --- \\
& Spectral FNO           & $\mathbf{0.959 \pm 0.053}$ & $\mathbf{0.167 \pm 0.010}$ & $\mathbf{0.135 \pm 0.006}$ & $+0.20$ \\
\midrule
\multirow{2}{*}{Heston-VG}  &
  Vanilla FNO            & $0.197 \pm 0.147$ & $0.391 \pm 0.232$ & $0.981 \pm 0.116$ & --- \\
& Spectral FNO           & $\mathbf{0.952 \pm 0.046}$ & $0.508 \pm 0.163$ & $2.110 \pm 1.797$ & $\mathbf{+0.76}$ \\
\midrule
\midrule
\multirow{2}{*}{BS-L\'evy}  &
  Vanilla FNO            & $0.367 \pm 0.073$ & $0.810 \pm 0.019$ & $0.485 \pm 0.149$ & --- \\
& Spectral FNO           & $\mathbf{0.815 \pm 0.165}$ & $\mathbf{0.117 \pm 0.011}$ & $\mathbf{0.175 \pm 0.029}$ & $+0.45$ \\
\midrule
\multirow{2}{*}{FPL}        &
  Vanilla FNO            & $0.526 \pm 0.306$ & $2.046 \pm 0.937$ & $1.119 \pm 0.513$ & --- \\
& Spectral FNO           & $\mathbf{0.967 \pm 0.055}$ & $\mathbf{0.090 \pm 0.041}$ & $\mathbf{0.077 \pm 0.023}$ & $+0.44$ \\
\bottomrule
\end{tabular}
\end{table}

\begin{figure}[t]
  \centering
  \includegraphics[width=0.85\linewidth]{fig16_spectral_fno_var}
  \caption{Variance ratio: Vanilla FNO vs.\ Spectral FNO across five
  problems. Target $= 1.0$ (dashed line). Spectral FNO closes the
  variance gap on every L\'evy-driven problem (Burgers-VG, Heston-VG,
  BS-L\'evy, FPL). See Figure~\ref{fig:spectral-fno-logpsd} for the
  corresponding spectrum comparison.}
  \label{fig:spectral-fno-var}
\end{figure}

\begin{figure}[t]
  \centering
  \includegraphics[width=\linewidth]{fig17_spectral_fno_logpsd}
  \caption{Spectrum matching: Vanilla FNO vs.\ Spectral FNO,
  $y$-axis on $\log_{10}$ scale. Left: log-spatial PSD MSE. Right:
  log-temporal PSD MSE. Dotted line: Euler--Mascheroni asymptotic
  floor ($0.063$). Spectral FNO improves spatial PSD on 3/4 problems
  (Heston-VG is the regression) and on FPL and Burgers-VG sits
  within a factor of two of the floor.
  Temporal PSD shows a trade-off on Heston-VG, an artefact of
  the optimiser allocating capacity to spatial structure and
  variance at temporal expense.}
  \label{fig:spectral-fno-logpsd}
\end{figure}

The pattern in Table~\ref{tab:spectral-fno} and
Figures~\ref{fig:spectral-fno-var}--\ref{fig:spectral-fno-logpsd}
is the expected one given the loss design.

\paragraph{Variance gap closure.}
On every L\'evy-driven problem the Spectral FNO closes the variance
gap to within $0.05$ of unity. The most dramatic case is Heston-VG,
where the gap shrinks from $0.80$ (vanilla) to $0.05$ (Spectral) ---
a sixteen-fold improvement. Burgers-VG, BS-L\'evy, and FPL show

exception: the vanilla FNO already saturates the variance metric
($0.993 \pm 0.006$), and the Spectral FNO regresses slightly
($0.948 \pm 0.019$). This regression is consistent with the
no-free-lunch behaviour visible elsewhere in the table: when there
is nothing to gain on the variance term, the optimiser spends the
gradient budget on the spectral terms and the variance suffers
marginally.

\paragraph{Spatial spectrum.}
Log-spatial PSD MSE drops by $11\times$ on Burgers-VG, $7\times$ on

Heston-VG it regresses ($0.39 \to 0.51$); the variance ratio gain is
paid for at the cost of both spatial and temporal spectrum quality.
On Burgers-VG, BS-L\'evy, and FPL the Spectral FNO sits within a
factor of two to three of the asymptotic floor of $0.063$
(Section~\ref{sec:theory}), suggesting near-optimal second-order
matching at the resolution we consider.

\paragraph{Temporal spectrum.}
Log-temporal PSD MSE shows the cleanest split. On Burgers-VG,
BS-L\'evy, and FPL the Spectral FNO improves the temporal spectrum
$\sim 9\times$, $\sim 3\times$, and $\sim 15\times$ respectively.

the Heston-VG regression to the structural CIR limit

no-free-lunch trade-off as the variance ratio, exacerbated here by
the periodogram of vorticity having a flatter and more demanding
temporal structure than the L\'evy benchmarks.

\paragraph{Pointwise accuracy.}
We report relative $L^2$ in the supplementary CSV. The Spectral FNO
sacrifices pointwise accuracy on every problem: relative $L^2$
worsens from $0.52 \to 0.76$ on Burgers-VG, $0.06 \to 1.53$ on

and $0.38 \to 1.08$ on FPL. The trade-off is significant: the
Spectral FNO is a \emph{distributional} surrogate, not a
\emph{trajectory} surrogate. Combining the spectral loss with a
small auxiliary $L^2$ penalty would plausibly recover pointwise
accuracy at the cost of partially reopening the variance gap; we
flag this as future work in Section~\ref{sec:discussion}.

\subsection{Reading the result}
\label{subsec:fno-reading}

The Spectral FNO results extend the central claim of the paper from
PINNs to neural operators. The same three-term loss that closes
the variance gap on a single-realisation PINN also closes it on a
many-trajectory operator surrogate, on four problems,
without any change to the architecture beyond the loss
augmentation. The exceptions are pedagogically useful:

variance metric, the spectral loss redistributes capacity rather
than improving the joint score; and Heston-VG shows that a problem
with structural difficulty
(Remark~\ref{rem:cir}, Section~\ref{sec:theory}) admits the same
local minimum even when the architecture changes from a PINN to an
FNO. The mechanism is the loss, not the architecture.

We do not claim that Spectral FNO supersedes vanilla FNO as a
trajectory surrogate. Reviewers seeking the lowest pointwise error
on a single SPDE realisation should still use the vanilla MSE
objective. The contribution here is narrower: an architecture-agnostic
loss recipe for closing the distributional gap on stochastic PDEs
across two distinct surrogate families (PINN and FNO).

%==============================================================================

%==============================================================================
\section{Conclusion}
\label{sec:conclusion}
\end{document}
